%%
%% fields.tex
%% 
%% Made by Alex Nelson <pqnelson@gmail.com>
%% Login   <alex@lisp>
%% 
%% Started on  2025-12-28T09:35:16-0800
%% Last update 2025-12-28T09:35:16-0800
%% 

\chapter{Field Theory}

\section{Fields}

\begin{definition}[Field]\mml{vectsp_1}
A \define{Field} is a commutative skew-field.
\end{definition}

\begin{theorem}[Kronecker's construction]\mml{field_1:def 3}
Let $\FF$ be field, let $p(x)\in\FF[x]$ be an irreducible polynomial.
Then $\FF[x]/(p(x))$ is a field.
\end{theorem}

\begin{proposition}[Well-definedness of characteristic of field]
For any field $\FF$, recall (\S\ref{ex:ring:canHomInt}) there exists a
unique [ring] morphism $\phi\colon\ZZ\to\FF$ such that $\phi(1)=1$
with $\ker(\phi)=(n)$ for some $n\geq0$.
Then there exists a unique induced morphism
$\widetilde{\phi}\colon\ZZ/n\ZZ\to\FF$ which is injective and either
$n=0$ or $n$ is a prime number.
\end{proposition}

\begin{definition}[Characteristic of a ring]\mml{ring_3:def 5}
Let $R$ be a ring. The \define{Characteristic} of $R$ is the
number $\Char(R)$ defined as:

If there exists an $n>0$ such that $(n)=\ker(\phi)$ where $\phi\colon\ZZ\to R$
sends $\phi(1)=1$, then $\Char(R)=n$ if it exists.

Otherwise, if for all $n\in\ZZ$ we have $n\cdot 1_{R}\neq0_{R}$, then
the characteristic of $R$ is defined to be zero $\Char(R)=0$.
\end{definition}

\begin{example}
The field $\FF_{p}$ has characteristic $p$, so does the ring $\FF_{p}[x]$.
\end{example}

\begin{example}
The rings $\ZZ$, $\QQ$, $\RR$, $\CC$ all have characteristic zero.
\end{example}

\begin{definition}[Prime Subfield]\mml{ring_3:def 10}
Let $\FF$ be a field. The \define{Prime Subfield} of $\FF$ is its
smallest subfield.
\end{definition}

\begin{proposition}[Prime subfield of field with characteristic zero is isomorphic to rationals]\mml{ring_3:103}
If $\FF$ is a field with $\Char(\FF)=0$, then its prime subfield is
isomorphic to $\QQ$.
\end{proposition}

\begin{proposition}[Prime subfield of finite characteristic field]\mml{ring_3:108}
Let $\FF$ be a field with $\Char(\FF)=p\neq0$, then its prime subfield
is isomorphic to the finite field with $p$ elements $\FF_{p}$.
\end{proposition}

\begin{example}
The prime subfield of $\FF_{p}(x)$ is $\FF_{p}$.
\end{example}

\section{Field Extension}

\begin{definition}[Field and ring extensions]\mml{field_4 field_4:7}
Let $\FF$ be a field. A \define{Field Extension} of $\FF$, denoted
$\kk/\FF$, is a field $\kk$ containing $\FF$ as a subfield $\FF\subset\kk$.
We say ``$\kk$ is \emph{over} $\FF$''.

More generally, if $R$ is a ring, we may define a \define{Ring Extension}
of $R$, denoted $S/R$, is a ring $S$ containing $R$ as a subring
$S\subset R$.
\end{definition}

\begin{example}
Some simple examples of field extensions $\QQ(\I)/\QQ$, $\CC/\RR$, and
$\FF_{p}(x)/\FF_{p}$.

However, $\FF_{p}[x]/\FF_{p}$ is not a field extension, but it is a
ring extension.
\end{example}

\begin{example}
The field $\QQ(\I)$ is the field obtained by \define{Adjoining} a root
of $x^{2}+1$ to $\QQ$.
\end{example}

\begin{proposition}[Polynomial of a ring is a polynomial of a ring extension]\mml{field_4:8}
Let $R$ be a ring, let $p(x)\in R[x]$ be a formal polynomial. For any
ring extension $S/R$, we see $p(x)\in S[x]$.
\end{proposition}

\begin{proposition}[Ring extension induces extension of formal polynomials]\mml{ring_4:19}
For any ring $R$ and ring extension $S/R$, the ring of formal
polynomials $S[x]$ is a ring extension of $R[x]$.
\end{proposition}

\begin{theorem}[Existence of field extension containing a root of a given polynomial]
Let $\FF$ be a field and $p(x)\in\FF[x]$ be a nonconstant, irreducible
polynomial. Then there always exists a field extension $\kk/\FF$ such
that $p$ has a root in $\kk$.
\end{theorem}

\begin{definition}[Polynomials splits in a ring]\mml{field_4:def 5}
Let $R$ be a ring, let $S$ be a nondegenerate ring (where $1_{S}\neq0_{S}$).
Let $p(x)\in R[x]$ be a polynomial. We say $p$ \define{Splits In} $S$ if
we can write $p(x)=a\prod^{n}_{i=1}(x-c_{i})$ where $a,c_{i}\in S$.
\end{definition}

\begin{definition}[Vector space underlying a field extension]\mml{field_4:def 6}
Let $\FF$ be a field, let $\kk/\FF$ be a field extension. Then we
define the \define{Vector Space} $V_{\kk}$ to be $\kk$ considered as a
vector space over $\FF$.
\end{definition}

\begin{definition}[Degree of field extension]\mml{field_4:def 7}
Let $\FF$ be a field, let $\kk/\FF$ be a field extension.
We define the \define{Degree} of $\kk/\FF$ to be the integer
$[\kk:\FF]=\dim_{\FF}(\kk)$ if $\kk$ is a finite-dimensional vector
space over $\FF$; otherwise, the conventions vary, but we will set
$[\kk:\FF]=\infty$. (Some references set $[\kk:\FF]=-1$ if $\kk$ is
not a finite-dimensional $\FF$-vector space, but I feel uneasy with
this convention.)

Moreover, when the degree of $\kk/\FF$ is finite, we say that it is a
\define{Finite Extension} of $\FF$.
\end{definition}